% !PRINTED 448
\weberpage{448}
If, namely, $x$ is a two-fold pole, and $x$ passes by $\omega$ into $x'$ and by $\Theta$ into $x''$, then not only are $x'$ and $x''$ different from $x$, but they are also different from one another, because, if they were identical, $x$ would remain unchanged by $\omega\Theta^3$. By $\omega\Theta$, $x$ then passes into a fourth position $x'''$, and $x'''$ is different both from $x'$ and from $x''$, because otherwise $x$ would remain unchanged by $\omega\Theta\omega = \Theta$ or by $\omega\Theta\Theta^{-1} = \omega$. Since $x$ remains unchanged by $\Theta^2$, $x$ passes by $\Theta^3$ into $x''$, by $\Theta^2\omega = \omega\Theta^2$ into $x'$, and by $\omega\Theta^3$ into $x'''$.

Since we have two points $c_1$, $c_2$ on each principal axis, there are 42 in all, and each of them can be carried into every other by substitutions of the group. From this it follows that these poles are not more than four-fold. Hence the theorem:

\noindent 4.\ \gesperrt{There are 21 eight-fold and 42 four-fold poles. On each principal axis lie four eight-fold and two four-fold poles. Through each eight-fold pole pass four principal axes. Every eight-fold pole can be carried into every other eight-fold pole, and every four-fold pole into every other four-fold pole.}

\begin{center}
\S.\,119.

\gesperrt{The three- and six-fold poles.}
\end{center}

We now turn to the consideration of the substitutions of the third order and of their poles, and choose as representative of such a substitution

\[
\chi = \begin{pmatrix} 0 & 0 & 1\\ 1 & 0 & 0\\ 0 & 1 & 0 \end{pmatrix},
\]
whose poles are obtained from

\[
\lambda\,x_1 = x_3, \quad \lambda\,x_2 = x_1, \quad \lambda\,x_3 = x_2
\]
From this it follows that $\lambda^3 = 1$; hence $\lambda$ is a third root of unity, and if therefore $\varrho$ denotes an imaginary third root of unity, the three poles result:

\weq{1}{\begin{array}{l}
x_1 = x_2 = x_3\\
x_1 = \varrho\,x_2 = \varrho^2 x_3\\
x_1 = \varrho^2 x_2 = \varrho\,x_3.
\end{array}}
